% main.tex

% https://tex.stackexchange.com/a/748667
\unprotect
    %%% Defined similarly to \obeyMPboxdepth, \ignoreMPboxdepth,
    %%% \obeyMPboxorigin, and \normalMPboxdepth
    \permanent\protected\def\strutMPboxdepth{%
        \let\meta_relocate_box\meta_strut_box_depth%
    }

    \protected\def\meta_strut_box_depth{%
        \setbox\b_meta_graphic\vpack{%
            \box\b_meta_graphic%
            %%% Add some depth to the box without affecting its size
            \nohrule height-\openstrutdepth width 0pt depth \openstrutdepth\relax%
        }
    }

    %%% The pagecolumns output routine alredy adds a \vfil elsewhere, so the one
    %%% here isn't needed and disrupts the depth added by \strutMPboxdepth.
    \let\page_otr_fill_page_unchecked\relax
\protect

%%% Comment the line below to confirm that an overfull vbox is generated;
%%% uncomment it to confirm that the overfull vbox message goes away.
\strutMPboxdepth
\setupbodyfont[pagella]


\definepapersize
  [youtubeHD]
  [width=1920bp,height=1080bp]
\setuppapersize[youtubeHD][youtubeHD]
\setuplayout[
    location=middle,
    width=middle,
    height=middle,
    topspace=0pt,
    backspace=0pt,
    header=0pt,
    footer=0pt
]
\defineframed[myframed][
    width=800pt,
    align=normal,
    offset=1cm,
    frame=off,
    indenting={yes,medium}
]
\setupbodyfont[modern,24pt]
\ctxloadluafile{parametric.lua}
\processMPfigurefile{register_mp_cmd.mp}

\startMPinclusions
def render_theorem_box =
numeric w, h;
w := CurrentWidth;
h := CurrentHeight;

path p;
p := (-w/2, -h/2) -- (w/2, -h/2) -- (w/2, h/2) -- (-w/2, h/2) -- cycle;
setbounds currentpicture to p;

picture theorem;
theorem := textext("\getbuffer[BufA]");

pair center_anchor, topleft_anchor;
center_anchor := center theorem;
topleft_anchor := ulcorner theorem;

enddef;


def jasper_thin =
    withpen pencircle scaled .1cm
enddef ;

def jasper_thick =
    withpen pencircle scaled .2cm
enddef ;

def jasper_lightgray =
    withcolor (0.6)
enddef ;

def jasper_lightgraypenthin =
    jasper_lightgray withpen pencircle scaled .1cm
enddef ;

def jasper_lightgraypenthick =
    jasper_lightgray withpen pencircle scaled .2cm
enddef ;


def jasper_darkgray =
    withcolor (0.4)
enddef ;

def jasper_darkgraypenthin =
    jasper_darkgray withpen pencircle scaled .15cm
enddef ;

def jasper_darkgraypenthick =
    jasper_darkgray withpen pencircle scaled .3cm
enddef ;
\stopMPinclusions
\starttext
    \startbuffer[BufA]
        \startframed[myframed]
            {\bf AXIOM 3} One and only one segment congruent to 
            a given segment can be laid off on a ray starting 
            from the origin.\par
            \color[white]{If two segments laid off on a ray from its vertex
            are non-congruent, their second endpoints do not 
            coincide, and the segment whose second endpoint 
            lies within the interior of the other segment is 
            said to be the least of the two. The other segment 
            is the greatest of the two.}
        \stopframed
    \stopbuffer
    % Display the theorem initially, in a prominent location.
    \dorecurse{7*30}{%
        \startMPpage
            render_theorem_box ;
            numeric t;
            t := 0;
            pair anchor;
            anchor := (1 - t) * center_anchor + t * topleft_anchor + t * (-2cm, -9cm);
            label(theorem, (0,3cm) - anchor);
        \stopMPpage%
    }
    % Translate the theorem to the top right corner.
    \dorecurse{2*30+1}{%
        \startMPpage
            render_theorem_box ;
            numeric t;
            t := (#1-1)/60;
            pair anchor;
            anchor := (1 - t) * center_anchor + t * topleft_anchor + t * (-2cm, -9cm);
            label(theorem, (0,3cm) - anchor);
        \stopMPpage%
    }
    \startMPinclusions
        def mythm =
            render_theorem_box ;
            numeric t;
            t := 1;
            pair anchor;
            anchor := (1 - t) * center_anchor + t * topleft_anchor + t * (-2cm, -9cm);
            label(theorem, (0,3cm) - anchor);
        enddef ;
        runscript(
            "
                jasper.math.cx = -15
                jasper.math.cy = 5
            "
        ) ;
    \stopMPinclusions
    % Pause once the theorem is in place.
    \dorecurse{4*30}{%
        \startMPpage
            mythm ;
        \stopMPpage%
    }
    % draw the vertex of a ray
    \dorecurse{3*30}{%
        \startMPpage
            mythm ;
            jasper_point[
                fx="cx"
                ,fy="cy"
                ,fz="0"
            ] ;
            jasper_render_segments ;
        \stopMPpage%
    }
    % extend a ray from the vertex
    \dorecurse{3*30}{%
        \startMPpage
            mythm ;
            jasper_point[
                fx="cx"
                ,fy="cy"
                ,fz="0"
            ] ;
            jasper_curve[
                ustart="0"
                ,ustop="1"
                ,usamples="2"
                ,fx="cx+13*u*#1/90"
                ,fy="cy"
                ,fz="0"
                ,drawcommands="jasper_thin"
            ] ;
            jasper_render_segments ;
        \stopMPpage%
    }
    % pause after extending a ray from the vertex
    \dorecurse{3*30}{%
        \startMPpage
            mythm ;
            jasper_point[
                fx="cx"
                ,fy="cy"
                ,fz="0"
            ] ;
            jasper_curve[
                ustart="0"
                ,ustop="1"
                ,usamples="2"
                ,fx="cx+13*u"
                ,fy="cy"
                ,fz="0"
                ,drawcommands="jasper_thin"
            ] ;
            jasper_render_segments ;
        \stopMPpage%
    }
    % extend a thick line segment
    \dorecurse{3*30}{%
        \startMPpage
            mythm ;
            jasper_point[
                fx="cx"
                ,fy="cy"
                ,fz="0"
            ] ;
            jasper_curve[
                ustart="0"
                ,ustop="1"
                ,usamples="2"
                ,fx="cx+13*u"
                ,fy="cy"
                ,fz="0"
                ,drawcommands="jasper_thin"
            ] ;
            jasper_point[
                fx="cx+9*#1/90"
                ,fy="cy"
                ,fz="0"
            ] ;
            jasper_curve[
                ustart="0"
                ,ustop="1"
                ,usamples="2"
                ,fx="cx+9*u*#1/90"
                ,fy="cy"
                ,fz="0"
                ,drawcommands="jasper_thick"
            ] ;
            jasper_render_segments ;
        \stopMPpage%
    }
    \startMPinclusions
        def rayseg =
            jasper_point[
                fx="cx"
                ,fy="cy"
                ,fz="0"
            ] ;
            jasper_curve[
                ustart="0"
                ,ustop="1"
                ,usamples="2"
                ,fx="cx+13*u"
                ,fy="cy"
                ,fz="0"
                ,drawcommands="jasper_thin"
            ] ;
            jasper_point[
                fx="cx+9"
                ,fy="cy"
                ,fz="0"
            ] ;
            jasper_curve[
                ustart="0"
                ,ustop="1"
                ,usamples="2"
                ,fx="cx+9*u"
                ,fy="cy"
                ,fz="0"
                ,drawcommands="jasper_thick"
            ] ;
        enddef ;
    \stopMPinclusions
    % pause after extending a thick line segment
    \dorecurse{3*30}{%
        \startMPpage
            mythm ;
            rayseg ;
            jasper_render_segments ;
        \stopMPpage%
    }
    \startbuffer[BufA]
        \startframed[myframed]
            {\bf AXIOM 3} One and only one segment congruent to 
            a given segment can be laid off on a ray starting 
            from the origin.\par
            If two segments laid off on a ray from its vertex
            are non-congruent, their second endpoints do not 
            coincide, \color[white]{and the segment whose second endpoint 
            lies within the interior of the other segment is 
            said to be the least of the two. The other segment 
            is the greatest of the two.}
        \stopframed
    \stopbuffer
    % If two segments laid off on a ray from its vertex
    % are non-congruent, their second endpoints do not 
    % coincide,
    \dorecurse{10*30}{%
        \startMPpage
            mythm ;
            rayseg ;
            jasper_render_segments ;
        \stopMPpage%
    }
    \startbuffer[BufA]
        \startframed[myframed]
            {\bf AXIOM 3} One and only one segment congruent to 
            a given segment can be laid off on a ray starting 
            from the origin.\par
            If two segments laid off on a ray from its vertex
            are non-congruent, their second endpoints do not 
            coincide, and the segment whose second endpoint 
            lies within the interior of the other segment is 
            said to be the least of the two. \color[white]{The other segment 
            is the greatest of the two.}
        \stopframed
    \stopbuffer
    % and the segment whose second endpoint 
    % lies within the interior of the other segment is 
    % said to be the least of the two.
    \dorecurse{13*30}{%
        \startMPpage
            mythm ;
            rayseg ;
            jasper_render_segments ;
        \stopMPpage%
    }
    \startbuffer[BufA]
        \startframed[myframed]
            {\bf AXIOM 3} One and only one segment congruent to 
            a given segment can be laid off on a ray starting 
            from the origin.\par
            If two segments laid off on a ray from its vertex
            are non-congruent, their second endpoints do not 
            coincide, and the segment whose second endpoint 
            lies within the interior of the other segment is 
            said to be the least of the two. The other segment 
            is the greatest of the two.
        \stopframed
    \stopbuffer
    % The other segment 
    % is the greatest of the two.
    \dorecurse{6*30}{%
        \startMPpage
            mythm ;
            rayseg ;
            jasper_render_segments ;
        \stopMPpage%
    }
    % add interior segment
    \dorecurse{4*30}{%
        \startMPpage
            mythm ;
            rayseg ;
            jasper_curve[
                ustart="0"
                ,ustop="1"
                ,usamples="2"
                ,fx="cx+6*cos(tau/8)*u"
                ,fy="cy+6*sin(tau/8)*u"
                ,fz="5"
                ,drawcommands="jasper_lightgraypenthick"
            ] ;
            jasper_point[
                fx="cx"
                ,fy="cy"
                ,fz="5"
                ,drawcommands="jasper_lightgray"
            ] ;
            jasper_point[
                fx="cx+6*cos(tau/8)"
                ,fy="cy+6*sin(tau/8)"
                ,fz="5"
                ,drawcommands="jasper_lightgray"
            ] ;
            jasper_render_segments ;
        \stopMPpage%
    }
    % pause
    \dorecurse{3*30}{%
        \startMPpage
            mythm ;
            rayseg ;
            jasper_curve[
                ustart="0"
                ,ustop="1"
                ,usamples="2"
                ,fx="cx+6*cos(tau/8)*u"
                ,fy="cy+6*sin(tau/8)*u"
                ,fz="5"
                ,drawcommands="jasper_lightgraypenthick"
            ] ;
            jasper_point[
                fx="cx"
                ,fy="cy"
                ,fz="5"
                ,drawcommands="jasper_lightgray"
            ] ;
            jasper_point[
                fx="cx+6*cos(tau/8)"
                ,fy="cy+6*sin(tau/8)"
                ,fz="5"
                ,drawcommands="jasper_lightgray"
            ] ;
            jasper_render_segments ;
        \stopMPpage%
    }
    % lay off on the ray
    \dorecurse{3*30}{%
        \startMPpage
            mythm ;
            rayseg ;
            jasper_curve[
                ustart="0"
                ,ustop="1"
                ,usamples="2"
                ,fx="cx+6*cos(tau/8-((#1/90)*(tau/8)))*u"
                ,fy="cy+6*sin(tau/8-((#1/90)*(tau/8)))*u"
                ,fz="5"
                ,drawcommands="jasper_lightgraypenthick"
            ] ;
            jasper_point[
                fx="cx"
                ,fy="cy"
                ,fz="5"
                ,drawcommands="jasper_lightgray"
            ] ;
            jasper_point[
                fx="cx+6*cos(tau/8-((#1/90)*(tau/8)))"
                ,fy="cy+6*sin(tau/8-((#1/90)*(tau/8)))"
                ,fz="5"
                ,drawcommands="jasper_lightgray"
            ] ;
            jasper_render_segments ;
        \stopMPpage%
    }
    \dorecurse{15*30}{%
        \startMPpage
            mythm ;
            rayseg ;
            jasper_curve[
                ustart="0"
                ,ustop="1"
                ,usamples="2"
                ,fx="cx+6*cos(0)*u"
                ,fy="cy+6*sin(0)*u"
                ,fz="5"
                ,drawcommands="jasper_lightgraypenthick"
            ] ;
            jasper_point[
                fx="cx"
                ,fy="cy"
                ,fz="5"
                ,drawcommands="jasper_lightgray"
            ] ;
            jasper_point[
                fx="cx+6*cos(0)"
                ,fy="cy+6*sin(0)"
                ,fz="5"
                ,drawcommands="jasper_lightgray"
            ] ;
            jasper_render_segments ;
        \stopMPpage%
    }
\stoptext 
